\documentclass[11pt]{article}

\usepackage[a4paper,top=0.7in,bottom=0.7in,left=0.85in,right=0.85in]{geometry}
\usepackage[T1]{fontenc}
\usepackage{lmodern}
\usepackage{microtype}
\usepackage{amsmath}
\usepackage{booktabs}
\usepackage{tabularx}
\usepackage{array}
\usepackage{enumitem}
\usepackage{xcolor}
\usepackage[hidelinks]{hyperref}
\usepackage{titlesec}

\titleformat{\section}{\large\bfseries}{\thesection}{0.6em}{}
\titleformat{\subsection}{\normalsize\bfseries}{\thesubsection}{0.5em}{}
\titlespacing*{\section}{0pt}{0.4\baselineskip}{0.1\baselineskip}
\titlespacing*{\subsection}{0pt}{0.3\baselineskip}{0.05\baselineskip}

\setlist{topsep=0pt,itemsep=0pt,parsep=0pt,leftmargin=1.1em}
\linespread{0.95}

\newcolumntype{Y}{>{\raggedright\arraybackslash}X}
\setlength{\emergencystretch}{3em}
\hbadness=10000

\title{\vspace{-1.2cm}Bitcoin Wallet Risk Analyzer\\[2pt]\large Project Status Report}
\author{Ori Sinvani \textnormal{(325770824)} \quad Harel Brener \textnormal{(214179012)} \quad Nehoray Chalfon \textnormal{(325833531)}\\[3pt]
\normalsize\normalfont Project Advisor: Ofer Zevin}
\date{May 27, 2026}

\begin{document}
\maketitle
\vspace{-1.0cm}

\noindent\textbf{1. Is the development complete?}\quad
Yes -- the full system is working end to end. The offline side covers everything from preparing
the data and building the wallet graphs, through training and calibrating the model, to producing
an evaluation report and an importance score for each input feature. The user-facing side is a
React web application backed by a FastAPI server that lets a visitor paste a Bitcoin address and
receive a verdict, an explanation, and supporting transaction information.

\noindent\textbf{2. What is missing?}\quad
The main gaps are polish rather than core functionality: there is no automated test suite yet;
the website is a single page and would benefit from a history view and a side-by-side comparison
view; and the configuration that allows the production website to talk to the production API is
still applied directly on the server rather than committed to the repository.

\noindent\textbf{3. Is the system deployed?}\quad
Yes. The website is hosted on the BGU CS department servers and is reachable at
\href{https://cryptotrace.cs.bgu.ac.il/}{\texttt{https://cryptotrace.cs.bgu.ac.il/}}. In
addition to the web interface, we ship a Python (Jupyter / Google Colab) notebook that loads the
trained model and walks through the same analysis on any address the reader supplies, so the
model can be demonstrated end-to-end without needing access to our server.

\noindent\textbf{4. Have you conducted experiments?}\quad
Yes. We compared three models head to head, all judged on the same wallets and the same
labels: (i) our main graph neural network, which reads both the wallet's own behaviour and the
pattern of its transactions with its neighbours; (ii) an XGBoost classifier that sees only the
wallet's own twelve summary statistics, with no neighbourhood information; and (iii) a simpler
graph baseline that sees the neighbourhood but ignores the details of each transaction. Comparing
the three lets us separate the contribution of each ingredient.

\noindent\textbf{5. How were the experiments conducted?}\quad
We combined two public datasets of labelled Bitcoin wallets, REAL-CATS (criminal wallets) and
Elliptic++ (a mix of criminal and benign), removed overlap between them, and balanced the result
so that exactly half of the wallets are criminal. This produced 86{,}870 wallets for training
and 21{,}718 held out for testing. For every wallet we asked the public \texttt{mempool.space}
service for its recent transactions and built a small graph around it: the wallet sits in the
middle, the parties it transacted with sit around it, and each link records how much was sent,
in which direction, and when. The main model was trained over multiple epochs with the usual
safeguards (a validation split, early stopping, and a calibration step so the predicted
probabilities can be read as honest confidence values). All three models were then evaluated on
the same held-out wallets so the comparison is fair, and we produced the standard battery of
diagnostic plots: confusion matrix, ROC curve, reliability diagram, and a per-feature importance
score.

\noindent\textbf{6. Important results.}\quad
On the held-out test set:

\vspace{-0.2em}
\begin{center}
\small
\begin{tabular}{lcccc}
\toprule
\textbf{Model} & \textbf{F\textsubscript{1}} & \textbf{Accuracy} & \textbf{Recall} & \textbf{ROC-AUC} \\
\midrule
Our graph model & \textbf{0.856} & \textbf{0.913} & \textbf{0.897} & \textbf{0.961} \\
XGBoost (wallet stats only) & 0.835 & 0.836 & 0.827 & 0.920 \\
Simple graph baseline & 0.633 & 0.761 & 0.674 & 0.809 \\
\bottomrule
\end{tabular}
\end{center}
\vspace{-0.2em}

Two gaps tell the story. The jump from the simple graph baseline to our model shows that the
details of each transaction (amount, direction, timing) matter a great deal. The jump from
XGBoost to our model shows that looking at the wallet's neighbourhood adds information that the
wallet's own statistics do not capture. The calibration check came back essentially flat, so the
confidence scores the website displays can be trusted at face value.

\noindent\textbf{7. What is still lacking, and the schedule going forward.}\quad
\textbf{Now -- mid-June:} add a held-out adversarial test set of recent wallets in neither
REAL-CATS nor Elliptic++ to stress-test generalisation.
\textbf{Mid-June -- end of June:} write the final report, polish the live demo, and build the
presentation deck. \textbf{Early July:} project defence and final submission.

\end{document}
